% Buku Ajar Kriptografi
% Berbasis silabus_kriptografi.tex dan RPS
% Program Studi Matematika S1
% Format standar buku ajar perguruan tinggi Indonesia

\documentclass[12pt,a4paper]{book}
% ============================================================
% PREAMBLE BUKU AJAR KRIPTOGRAFI
% ============================================================

% ------------------------------------------------------------
% Encoding & Bahasa
% ------------------------------------------------------------
\usepackage[utf8]{inputenc}
\usepackage[T1]{fontenc}
\usepackage[indonesian]{babel}

% ------------------------------------------------------------
% Layout & Tipografi
% ------------------------------------------------------------
\usepackage[verbose=false]{geometry}
\geometry{
    a4paper,
    top=2.5cm,
    left=3cm,
    right=2cm,
    bottom=2cm
}

\setlength{\parindent}{0pt}
\setlength{\parskip}{0.5em}
\setlength{\emergencystretch}{2em}
\renewcommand{\baselinestretch}{1.5}
\raggedbottom
\usepackage{ragged2e}
\justifying

% ------------------------------------------------------------
% Paket Tabel & Struktur Dokumen
% ------------------------------------------------------------
\usepackage{booktabs}
\usepackage{longtable}
\usepackage{array}
\usepackage{enumitem}
\usepackage{subfiles}

% Nama job main document (untuk deteksi kompilasi per bab vs buku utuh)
\makeatletter
\def\MainDocJobname{buku_kriptografi}
\makeatother

% ------------------------------------------------------------
% Grafik, Diagram, & Hyperlink
% ------------------------------------------------------------
\usepackage{url}
\usepackage{graphicx}
\usepackage{xcolor}
\usepackage{tikz}
\usetikzlibrary{shapes.geometric, arrows.meta, positioning, calc}
\usepackage[hypertexnames=false]{hyperref}
\usepackage{bookmark}

% ------------------------------------------------------------
% Header & Footer
% ------------------------------------------------------------
\usepackage{fancyhdr}

% ------------------------------------------------------------
% Pengaturan Section (judul rata kiri, isi justified)
% ------------------------------------------------------------
\usepackage{titlesec}

\titleformat{\section}
  {\normalfont\Large\bfseries\raggedright}
  {\thesection}{1em}{}

\titleformat{\subsection}
  {\normalfont\large\bfseries\raggedright}
  {\thesubsection}{1em}{}

\titleformat{\subsubsection}
  {\normalfont\normalsize\bfseries\raggedright}
  {\thesubsubsection}{1em}{}

\titlespacing*{\section}{0pt}{3.5ex plus 1ex minus .2ex}{2.3ex plus .2ex}

% Pastikan isi paragraf tetap justified setelah judul
\usepackage{etoolbox}
\pretocmd{\section}{\justifying}{}{}
\pretocmd{\subsection}{\justifying}{}{}
\pretocmd{\subsubsection}{\justifying}{}{}

% ------------------------------------------------------------
% LISTING (Kode Program)
% ------------------------------------------------------------
\usepackage{listings}

\lstset{
    basicstyle=\ttfamily\small,
    breaklines=true,
    frame=single,
    numbers=left,
    numberstyle=\tiny,
    keywordstyle=\color{blue},
    commentstyle=\color{green!60!black},
    stringstyle=\color{red},
    showstringspaces=false,
    tabsize=2,
    float=false
}

\lstdefinestyle{matlab}{
    language=Matlab,
    morekeywords={function,end,if,else,for,while,return}
}

\lstdefinestyle{cstyle}{
    language=C,
    morekeywords={int,char,void,return,for,while,if,else}
}

\lstdefinestyle{cppstyle}{
    language=C++,
    morekeywords={int,char,void,return,for,while,if,else,string,cout,cin}
}

\lstdefinestyle{pascalstyle}{
    language=Pascal,
    morekeywords={program,var,begin,end,function,procedure,if,then,else,for,to,do,while}
}

% ------------------------------------------------------------
% KONTROL FLOAT (Gambar, Tabel, Listing)
% ------------------------------------------------------------
\usepackage{float}      % opsi [H]
\usepackage{placeins}   % \FloatBarrier

% Parameter float (lebih stabil untuk buku ajar)
\renewcommand{\topfraction}{0.9}
\renewcommand{\bottomfraction}{0.8}
\renewcommand{\textfraction}{0.05}
\renewcommand{\floatpagefraction}{0.85}
\setcounter{topnumber}{3}
\setcounter{bottomnumber}{2}
\setcounter{totalnumber}{4}

% Default posisi figure & table
\makeatletter
\def\fps@figure{!htbp}
\def\fps@table{!htbp}
\makeatother

% Cegah float melewati section
\let\oldsection\section
\renewcommand{\section}{\FloatBarrier\oldsection}

\let\oldsubsection\subsection
\renewcommand{\subsection}{\FloatBarrier\oldsubsection}

% Pastikan float selesai sebelum listing
\pretocmd{\lstlisting}{\FloatBarrier}{}{}

% ------------------------------------------------------------
% Metadata Dokumen
% ------------------------------------------------------------
\title{\textbf{Buku Ajar Kriptografi}\\
\large Untuk Program Studi Matematika S1}
\author{Departemen Matematika}
\date{}


\begin{document}
\newcommand{\InMainDocument}{}
\hypersetup{pageanchor=false}

%==============================================================================
% HALAMAN JUDUL LUAR (COVER)
%==============================================================================
\begin{titlepage}
\centering
\vspace*{2cm}
{\LARGE \textbf{BUKU AJAR}}\\[1cm]
{\Huge \textbf{KRIPTOGRAFI}}\\[0.5cm]
{\Large Untuk Program Studi Matematika S1}\\[0.8cm]
{\normalsize Mata kuliah Kriptografi \quad|\quad Kode MAT-XXX \quad|\quad 3 SKS}\\[0.3cm]
{\normalsize Program Studi Matematika \quad|\quad Fakultas MIPA}\\[0.3cm]
{\normalsize ISBN: \textit{(diisi sesuai kebijakan institusi)}}\\[1.6cm]
{\large Berbasis Outcome-Based Education (OBE)}\\[3cm]
{\large Departemen Matematika}\\[0.5cm]
{\large \today}
\end{titlepage}

%==============================================================================
% HALAMAN JUDUL DALAM
%==============================================================================
\newpage
\thispagestyle{empty}
\centering
\vspace*{2.5cm}
{\Huge \textbf{Buku Ajar Kriptografi}}\\[1cm]
{\Large Untuk Program Studi Matematika S1}\\[1.2cm]
{\large Berbasis Rencana Pembelajaran Semester (RPS)}\\[0.8cm]
{\large 14 Bab \quad|\quad 3 SKS}\\[1.2cm]
\vspace{0.5cm}
{\normalsize
Mata kuliah Kriptografi \quad|\quad Kode MAT-XXX \quad|\quad 3 SKS\\
Program Studi Matematika \quad|\quad Fakultas MIPA\\
Penerbit: \textit{Nama Institusi}, \textit{Kota}, \textit{Tahun}
}\\[2.5cm]
{\large Departemen Matematika}

%==============================================================================
% HALAMAN HAK CIPTA / IDENTITAS PENULIS (format BAN-PT/LAM)
%==============================================================================
\newpage
\thispagestyle{empty}
\vspace*{1.5cm}
\noindent
\textbf{Hak Cipta}\\[0.3em]
Hak cipta dilindungi undang-undang. Dilarang mengutip atau memperbanyak sebagian maupun seluruh isi buku ini tanpa izin tertulis dari penerbit.\\[1.2cm]

\noindent
\textbf{Penulis}\\[0.3em]
\textit{Nama Penulis}\\
Departemen Matematika, \textit{Nama Fakultas/Universitas}\\[0.8cm]

\noindent
\textbf{Penelaah}\\[0.3em]
\textit{Nama Penelaah}\\
\textit{Afiliasi}\\[0.8cm]

\noindent
\textbf{Editor}\\[0.3em]
\textit{Nama Editor} \quad (atau --- jika tidak ada)\\[0.8cm]

\noindent
\textbf{Penerbit}\\[0.3em]
\textit{Nama Institusi/Unit}, \textit{Kota}, \textit{Tahun}\\[0.5cm]

\noindent
\textbf{Tahun Terbit / Copyright}\\[0.3em]
\textit{Tahun}

%==============================================================================
% FRONTMATTER
%==============================================================================
\frontmatter
\newpage
\thispagestyle{empty}

% Kata Pengantar
\chapter*{Kata Pengantar}
\addcontentsline{toc}{chapter}{Kata Pengantar}

Puji syukur ke hadirat Tuhan Yang Maha Esa atas terselesaikannya Buku Ajar Kriptografi ini. Buku ini disusun sebagai bahan ajar mata kuliah Kriptografi untuk Program Studi Matematika jenjang Strata 1 (S1), berdasarkan Silabus dan Rencana Pembelajaran Semester (RPS) mata kuliah Kriptografi serta Capaian Pembelajaran Mata Kuliah (CPMK) Program Studi Matematika yang berlaku.

Materi dalam buku ini disusun mengikuti alur pembelajaran semester dengan 14 bab yang sesuai dengan 14 pertemuan efektif (tidak termasuk Ujian Tengah Semester dan Ujian Akhir Semester). Setiap bab mencakup tujuan pembelajaran, penyajian materi, dan rangkuman. Implementasi algoritma kriptografi disajikan dalam empat bahasa pemrograman: MATLAB, C, C++, dan Pascal, sesuai dengan tuntutan kurikulum.

Penulis berharap buku ajar ini dapat membantu mahasiswa dalam mempelajari kriptografi secara sistematis, baik dari aspek matematika maupun implementasi praktis. Saran dan masukan untuk perbaikan edisi selanjutnya sangat diharapkan.

\vspace{1cm}
\noindent\hfill Penulis

% Daftar Isi
\newpage
\tableofcontents

% Daftar Gambar
\newpage
\listoffigures
\addcontentsline{toc}{chapter}{\listfigurename}

% Daftar Tabel
\newpage
\listoftables
\addcontentsline{toc}{chapter}{\listtablename}

%==============================================================================
% MAINMATTER - 14 Bab sesuai RPS
%==============================================================================
\mainmatter
\hypersetup{pageanchor=true}

\subfile{chapters/chap01_pengantar}
\subfile{chapters/chap02_aritmetika}
\subfile{chapters/chap03_euclid}
\subfile{chapters/chap04_teori_bilangan}
\subfile{chapters/chap05_caesar}
\subfile{chapters/chap06_vigenere}
\subfile{chapters/chap07_block_cipher}
\subfile{chapters/chap08_rsa}
\subfile{chapters/chap09_hash}
\subfile{chapters/chap10_tanda_tangan}
\subfile{chapters/chap11_protokol}
\subfile{chapters/chap12_manajemen_kunci}
\subfile{chapters/chap13_kriptanalisis}
\subfile{chapters/chap14_implementasi}

%==============================================================================
% BACKMATTER
%==============================================================================
\backmatter

\chapter*{Daftar Referensi}
\addcontentsline{toc}{chapter}{Daftar Referensi}

\bibliographystyle{unsrt}
\bibliography{referensi}

\end{document}
